% Seccion: Planteamiento del problema
\section{Planteamiento del problema}
\label{sec:cap01_planteamiento}

Los colectores solares de aire convencionales presentan una eficiencia térmica inherentemente baja debido a las propiedades de transporte deficientes del aire.
El bajo número de Prandtl del aire limita el coeficiente de transferencia de calor por convección entre la placa absorbedora y la corriente fluida.
Esto se traduce en la formación de una capa límite térmica de gran espesor que actúa como una resistencia al paso de calor.
En la región de Río Cuarto, Córdoba, el secado solar de productos agrícolas es vital, pero requiere sistemas térmicos más eficientes.
Los diseños tradicionales de placa plana no logran satisfacer los caudales y temperaturas requeridas sin incurrir en grandes áreas de colección.
Por lo tanto, surge el problema de optimizar la geometría interna del absorbedor para inducir turbulencia localizada sin elevar excesivamente la caída de presión.
El diseño de un absorbedor tipo serpentín se presenta como una alternativa para maximizar el camino óptico del aire y romper de forma continua la subcapa límite laminar.
